\chapter*{Conclusion and Future Research}
\addcontentsline{toc}{chapter}{Conclusion and Future Research}
\markboth{Conclusion and Future Research}{Conclusion and Future Research}

\section*{What the project has completed}

This thesis converts eight specialised studies into one traceable research
architecture. Its main contribution is the separation and subsequent
connection of five objects: the historical gold series, the option-implied
gold-price law, production, unit cost and corporate cash flow. The resulting
chain is explicit enough to identify which evidence supports each output and
where a change of data, units, measure or accounting perimeter occurs.

Several tasks described as future work in the standalone articles are now
complete. The Team 8 project no longer stops at Black--Scholes, Heston or
independent Bates jumps: it contains the full affine Bates--Hawkes
specification, constrained calibration, residual and smile diagnostics,
four-model path simulation and a dense-only rolling comparison over
\OosDates{} origins and \OosObservations{} common forecasts. Its price
scenarios have also been connected, on a dated valuation clock, to the Team~4
forecast from the 2026Q3 stub and the unified DCF/WACC contract; Barrick's
Q1--Q2 2026 actuals precede the valuation time and serve only as level checks.
The Team~4 price demonstrator is excluded. The numerical integration is an
implemented four-branch scenario experiment; the physical forecasting and
corporate valuation stages remain incomplete.

The evidence was rebuilt on 24 October 2026. GLD's 0.40\% sponsor fee now
enters the calibration as carry; the valuation counts only flows after
2 September 2026; gold paths start from the XAUUSD spot of USD~4,391.32/oz at
the valuation time rather than from Barrick's earlier Q2 average realised
price; the daily path engines were corrected; Bates and Bates--Hawkes are
refined from their nested models; and the rolling test gives every model the
same optimizer search. Each correction moves the medians by less than the
numerical range reported below (Table~\ref{tab:unified-attribution}).

No per-share value is reported and no comparison with the stock quote is made: without a reconciled equity bridge such a comparison could not identify fair value, mispricing or physical probabilities. The copper extension likewise reports conditional aggregate sensitivities, about USD~\CuDeltaLow--\CuDeltaHigh~billion of explicit flows to 2030, rather than a validated company value.

The comparison further shows that, with production, cost and WACC shocks
held fixed, the choice of gold-price law hardly moves the central operating
proxy. Medians range from USD~\ValMedianLow{} to \ValMedianHigh~billion across
the four engines, whereas doubling the paths or changing the seed moves
single-model medians across USD~\NumericalMedianLow--\NumericalMedianHigh~billion:
the spread between models is within numerical noise and does not rank them.
The terminal value carries about \TerminalShare\% of the mean, so the
result depends more on the terminal convention than on the price model. This
is a useful attribution result: it identifies where price-model complexity
propagates through the DCF and where downstream assumptions compress or
dominate its effect.

\section*{Present limitations}

The completed integration remains conditional at several economic interfaces.
First, GLD option dynamics are estimated under \(\mathbb Q\), whereas a
long-horizon operating valuation requires scenarios justified under
\(\mathbb P\). Starting the option-implied engines at the XAUUSD spot
transfers distributional shape; it is not a validated GLD-to-gold conversion
or change of measure. The gold spot, Treasury curve, calibration surface and
valuation time are synchronous and dated 2 September 2026.

Second, the unified run excludes the elapsed 2026Q1--Q2 actuals and fixes the
Team~4 forecast from the 2026Q3 stub to 2030Q4. That design isolates price-engine risk but does not propagate reserve,
grade, recovery, energy, labour, mine-mix or closure uncertainty. Dependence
between gold prices and operating conditions is therefore not estimated.

Third, Team 5 supplies a coherent DCF architecture rather than a
filing-reconciled Barrick model. The operating-to-EBIT bridge, cash taxes,
working capital, sustaining and growth capex, debt, cash, minorities, streams,
other claims, non-operating assets and diluted shares require one dated
corporate reconciliation. No equity bridge is applied, so the outputs remain
aggregate operating proxies. Finite-mine closure also remains less developed than the
generic going-concern terminal block.

Finally, a sophisticated stochastic law does not eliminate model risk. The
nested calibration guarantees that each richer model fits the current
cross-section at least as well as the model it contains, so the in-sample
ordering is not evidence of superiority. In the rolling repricing test Full
Bates--Hawkes has the lowest date-equal IV RMSE, \OosWRMSE~bp, but its
advantage over Bates is not significant and no model beats interpolated IV
persistence. Its branching ratio is \HawkesBranching, below
unity; the Feller constraint binds for the jump models and mean reversion is
weakly identified. Parameter pressure,
surface sampling, optimizer dependence, regime shifts and the choice of
physical-measure transformation remain sources of uncertainty.

Costs retain a persistent trend without additional technology-driven reductions.
This is prudential relative to otherwise identical efficiency-improvement
scenarios, not a mathematical lower bound. Stochastic operating costs would
introduce favourable and adverse outcomes, and their conditional mean could
remain smooth. Productivity improvements and their investment costs must be
evaluated jointly with the growth already included in the DCF.

\section*{Future research programme}

The next research cycle must establish physical commodity forecasts and economic reconciliation. A dated Barrick data contract should
map mine-level production, realised prices, royalties, cash costs, AISC,
depreciation, sustaining and growth capex, working capital, taxes and closure
liabilities into EBIT, invested capital and FCFF. Debt, cash, minority
interests, streams, non-operating assets and diluted shares should then close
the EV-to-equity bridge to the same filing and listing.

The commodity interface requires a separate empirical programme. The GLD-to-
physical-gold basis should be estimated and stress-tested, and risk-neutral
option scenarios should be converted into physical scenarios through an
explicitly validated risk-premium or filtering model. Only after that step can
the long-horizon paths be interpreted as more than option-implied conditional
shapes.

Operating uncertainty should then be reintroduced jointly. Team 4 production
and cost models can be sampled alongside price, but their dependence structure
must be estimated or scenario-defined. Mine depletion, expansion projects,
reserve conversion, inflation, energy and foreign-exchange shocks should feed
an explicit finite-resource closure. The comparison can then be repeated to
measure whether gold-model choice or operating/accounting uncertainty explains
more of the equity-value distribution.

A later information layer may augment the filtration with timestamped
macroeconomic and geopolitical data. News embeddings, NLP classifications or
machine-learning forecasts are acceptable only with event-time alignment,
leakage controls, pre-declared benchmarks and true chronological evaluation.
The relevant question is not whether a flexible model improves in-sample fit,
but whether incremental information improves forecasts beyond price,
volatility and persistence benchmarks.

Further model research may examine time-varying risk premia, regime-dependent
parameters, model averaging and a rough-Hawkes option-pricing construction.
Those extensions should be judged through the same hierarchy used here:
mathematical admissibility, deterministic implementation tests, dated
out-of-sample evidence and economic propagation through a reconciled corporate
model.

\section*{Concluding synthesis}

The project demonstrates the value of interdisciplinarity in applied finance.
Probability, stochastic calculus, econometrics, numerical methods,
optimization, accounting and corporate valuation are not parallel subjects in
this setting; each constrains a different interface of the same experiment.
The strongest outcome is therefore not one point estimate or one preferred
model. It is a reproducible framework that keeps historical evidence,
option-implied information, operating assumptions and accounting claims
separate until the moment they are deliberately connected.

Within that framework, Full Bates--Hawkes offers the richest accepted
description of current option-implied variance and clustered jumps, while the
four-model valuation shows how little a better surface fit alone can establish
about a mining company's operating value. The thesis closes with a conditional result,
an explicit inventory of unresolved economic bridges and a research programme
capable of testing them. The copper extension of
Section~\ref{sec:copper-extension} expands the operating perimeter and adds
dated conditional surface tests, while leaving the accounting, realised-price,
mine-life and equity bridges explicit. Its aggregate paired values remain
operating proxies, not a corporate target.
The study does not issue a fair-value certification, target
price or investment recommendation.

% Physical forecasting gates are defined in Chapter 3.
